% =====================================================================
% Calibration methodology write-up — aligned with the committed model
% (calibrate.py: two-stage surrogate-assisted SMM; five-component
%  stylised-fact loss D = ΔKS + ΔV + ΔACF1 + ΔACF2 + ΔHill; Franke–Westerhoff
%  inverse-SD weighting via a Künsch moving-block bootstrap).
%
% SCOPE OF THIS DRAFT (per request): the story + the setup + the loss and its
% weighting only. The two optimisers (exhaustive grid; XGBoost surrogate) and
% the validation protocol (fresh-seed re-rank, moment coverage / J-test,
% ablation, robustness) come AFTER this and are NOT written here.
%
% Destination: Chapter 4 (Data, calibration and validation), opening sections.
% Preamble-less snippet; natbib author–year (\citep/\citet/\citealp), matching
% clearing_writeups.tex. Section labels reuse the existing calibration.tex names
% so cross-refs resolve; this draft refines those opening sections.
%
% Citation keys used (all resolve in introduction_refs.bib / model_refs.bib):
%   franke_westerhoff_2012, lamperti_2018, gao_2022_xgb, gao_2022_hfabm,
%   gao_2023, cont_2001, hill_1975, kunsch_1989, cont_stoikov_talreja_2008,
%   majewski_2018, simudyne_odd.
% Optional lineage keys to ADD (BibTeX at end of file): winker_gilli_maringer_2007,
%   grazzini_richiardi_2015.
% =====================================================================


% ============== §0 — CALIBRATION STRATEGY ============================

\section{Calibration strategy}
\label{sec:cal-strategy}

An agent-based market admits no tractable likelihood. The map from the
behavioural parameters to the realised price path runs through millions of
discrete order-book events, and there is no closed form for the density of the
data given the parameters, so the maximum-likelihood and Bayesian machinery of
classical estimation simply does not apply. The established response, and the
one taken here, is to match \emph{moments} rather than likelihoods: simulate the
model, summarise the simulated and the real market by the same battery of
statistics, and choose the parameters that make the two batteries as close as
possible. For financial-market agent-based models this \emph{simulated method of
moments} was put on a systematic footing by \citet{franke_westerhoff_2012},
whose structural stochastic-volatility models are fitted to the stylised facts
of returns by exactly this device; the loss construction below is taken from
their work (Section~\ref{sec:cal-loss}).

Where this calibration departs from the classical method is in how the loss is
\emph{minimised}. A single evaluation of the objective costs weeks of simulated
trading replicated over several random seeds, so the affordable budget is a few
hundred evaluations --- far too few to search a parameter space directly. The
remedy is a machine-learning \emph{surrogate}: a cheap statistical emulator of
the loss surface, trained on a modest design of true evaluations and then
optimised in place of the simulator. Surrogate-assisted calibration of
agent-based models was introduced by \citet{lamperti_2018} and developed, for
Chiarella-type models fitted to the same stylised-fact battery used here, by
\citet{gao_2022_xgb,gao_2022_hfabm}; the two-stage workflow --- a surrogate
search followed by a refinement on the true simulator --- and the exhaustive
grid that cross-checks it \citep{gao_2023} are described after the loss.

This separation organises the rest of the chapter. The \emph{target} and the
\emph{loss} --- common to every optimiser --- are fixed first (this section and
the next two); the two optimisers and the validation that guards them follow.
Two principles run throughout. The first is parsimony: the calibrated set is
kept deliberately small, anything that can be measured or pinned from outside
the model is removed from the search, and a richer specification earns its extra
parameters only if it improves the fit out of sample. The second is honesty
about identification: the loss surface is reported in full, and where a parameter
is weakly identified or rests against a search bound this is disclosed rather
than smoothed over.


% ============== §1 — DATA AND REGIMES ================================

\section{Data and regimes}
\label{sec:cal-data}

The empirical basis is the E-mini S\&P~500 futures tape: one-minute bars of the
front-month contract, rolled at volume crossover and restricted to regular
trading hours, with the quote midpoint $(b_t + a_t)/2$ as the price series. The
mid is the same observable the simulator produces, so every model--data
comparison is like-for-like. Returns are intraday log-returns of the mid;
overnight (session-boundary) returns are excluded on both sides, with the
session boundaries read from the data rather than imposed on a fixed grid ---
ES sessions vary in length, and a fixed-grid convention was found to leak
overnight gaps into the intraday moments.

Two regimes are calibrated as fully independent problems. The \emph{calm} regime
is drawn from a quiet stretch of 2019; the \emph{stressed} regime is the
COVID-19 crash window of late February to early April 2020. The two windows are
matched in length (about seventy-five sessions each) so that the
distribution-distance moment below, whose sampling scale depends on sample size,
is comparable across regimes. Calibrating the regimes separately --- with no
shared behavioural parameters --- is deliberate: the two market states differ in
volatility, book depth and microstructure, and forcing common parameters would
trade fit for compromise; the calm/stressed split follows the convention of the
reference model \citep{simudyne_odd}. A separate market-by-price (MBP-10) sample
of order-book events supports the direct measurement of order-placement depth.


% ============== §2 — WHAT IS CALIBRATED ==============================

\section{What is calibrated}
\label{sec:cal-split}

Parameters enter the model through three doors, and only one of them involves
the optimiser.

\emph{Measured offline.} Anything that can be estimated directly from data, is.
The fundamental-value process and its local volatility are fitted to the tape
(Chapter~3.\ref{sec:model-fv}); the order-placement depth $p_{\mathrm{zi}}$ is
estimated by maximum likelihood on the distance-from-mid distribution of new
limit orders in the MBP-10 data; and the loss-absorbing capital of members and
clients is taken from public FCM financials and CCP disclosures
(Chapter~3.\ref{sec:model-margin}). None of these re-enters the calibration
loop, with one disclosed exception noted below.

\emph{Pinned by design.} Structural choices follow their sources rather than the
optimiser: the momentum horizon $\lambda$ is fixed externally as in
\citet{majewski_2018}; the zero-intelligence market-order rate keeps the
reference model's baseline \citep{simudyne_odd}, itself structured after
\citet{cont_stoikov_talreja_2008}; and all clearing-tier constants are regulatory
pins (Chapter~3.\ref{sec:model-margin}). Pinning is justified case by case in the
ablation study --- a parameter is pinned only where calibrating it was shown to
change the fit negligibly.

\emph{Calibrated.} What remains is a small, regime-specific behavioural set. In
the calm regime the calibrated vector is three-dimensional,
$\theta = (c_f,\ \alpha_{\mathrm{zi}},\ \delta_{\mathrm{zi}})$: the
fundamental-trader belief-width scale and the zero-intelligence limit-order
arrival and cancellation rates. Under stress the same set is augmented with the
placement depth $p_{\mathrm{zi}}$ --- the one measured parameter that re-enters
the search --- giving four dimensions, because a sparser book materially improves
the stressed fit while the calm fit prefers the directly measured value. Search
bounds are literature-anchored: the zero-intelligence rates are confined to
ranges of the same order as their \citet{cont_stoikov_talreja_2008} baselines
(an early unconstrained run drove the market-order rate to many times its
reference value, producing a book-sweeping artefact), and the belief-width box
was tightened to the economically sensible basin located by diagnostic runs.


% ============== §3 — TARGETS: THE STYLISED-FACT BATTERY ==============

\section{Targets: the stylised-fact battery}
\label{sec:cal-targets}

The calibration targets are the stylised facts of high-frequency returns
\citep{cont_2001} --- the statistical fingerprints any credible simulated market
must reproduce. Three carry the weight. Returns are \emph{fat-tailed}: large
moves occur far more often than a Gaussian allows. Returns are \emph{linearly
unpredictable}: their autocorrelation is near zero from the first lag, so the
price is close to a martingale. And volatility \emph{clusters}: the
autocorrelation of \emph{absolute} returns is positive and decays slowly, so
turbulent minutes follow turbulent minutes. These facts are encoded in five
moment groups, computed identically on simulated and empirical intraday mid
returns.

The first group, \textbf{KS}, is the two-sample Kolmogorov--Smirnov statistic
$\sup_x |\widehat F_{\mathrm{sim}}(x) - \widehat F_{\mathrm{emp}}(x)|$ between the
simulated and empirical return distributions --- the largest vertical gap
between their cumulative curves, and the whole-distribution target of
\citet{gao_2022_xgb}; its goal is zero, a simulated distribution indistinguishable
from the real one. The second, \textbf{V}, is the return standard deviation,
which anchors the overall volatility scale \citep{gao_2022_hfabm}. The third,
\textbf{ACF1}, is the autocorrelation of returns at lags $\{1,5,10,20\}$ minutes,
each smoothed as a forward three-lag average (the mean of lags $\{c,c{+}1,c{+}2\}$,
following \citealp{gao_2022_xgb}); its target is the empirical near-zero profile,
so the group penalises any mechanical predictability the agents introduce, such
as bid--ask bounce or trend-chasing. The fourth, \textbf{ACF2}, is the
autocorrelation of \emph{absolute} returns at the same lags --- the
volatility-clustering target. Absolute returns are used in preference to squared
returns deliberately: $r^2$ is dominated by a handful of extreme observations,
which makes its autocorrelation so noisy that a clustering miss would vanish
under the standardisation of Section~\ref{sec:cal-loss}, whereas $|r|$ carries
the same signal with far less sampling noise \citep{cont_2001}. The fifth,
\textbf{Hill}, is the \citet{hill_1975} tail index of $|r|$, averaged over a band
of upper-tail fractions for stability; the index $\alpha$ in
$\Pr(|r|>x)\sim x^{-\alpha}$ summarises how fast the tail decays, with empirical
equity-index returns near the ``cubic law'' $\alpha\approx 3$ and smaller values
meaning fatter tails. Following \citet{gao_2022_hfabm}, Hill rather than kurtosis
is the tail target: excess kurtosis is governed by the few largest observations,
so matching it amounts to chasing noise, and it is tracked only as a diagnostic.

KS and Hill are paired on purpose, because each controls a failure mode the
other cannot. With KS alone the optimiser can match the body of the distribution
while the tail explodes; with Hill alone the tail is pinned but the body drifts.
Together they constrain the return distribution everywhere --- body and tail at
once.


% ============== §4 — THE LOSS FUNCTION ===============================

\section{The loss function}
\label{sec:cal-loss}

The five groups mix units --- a standard deviation of order $10^{-4}$, a
dimensionless KS statistic, a set of autocorrelations, a tail index near three
--- so they cannot simply be added. Following the inverse-variability weighting
of \citet{franke_westerhoff_2012}, each moment's distance is expressed in units
of that moment's own \emph{sampling noise} before the groups are combined:
\begin{equation}
\label{eq:cal-loss}
D(\theta)
\;=\!\!\sum_{c \in \{\mathrm{KS,\,V,\,ACF1,\,ACF2,\,Hill}\}}\!\! \Delta_c(\theta),
\qquad
\Delta_c(\theta)
\;=\; \frac{1}{|M_c|}\sum_{m \in M_c}
\frac{\bigl|\,\widehat m_{\mathrm{sim}}(\theta) - \widehat m_{\mathrm{emp}}\,\bigr|}{s_m}\,,
\end{equation}
where $M_c$ collects the moments in group $c$ and $s_m$ is the empirical
sampling standard deviation of moment $m$. Each group distance $\Delta_c$ is the
mean standardised miss over its moments, and the total loss is their sum, so the
five facts carry equal weight and $D$ is a single dimensionless number,
comparable across model versions.

\subsection{Weighting: noise standardisation}
\label{sec:cal-weight}

The weights $s_m$ make the heterogeneous moments commensurate. Each is the
standard deviation of moment $m$ under historical sampling alone, estimated once
per regime by a Künsch moving-block bootstrap \citep{kunsch_1989} of the
empirical one-minute mid returns: the series is resampled in contiguous blocks,
the whole moment battery is recomputed on each resample, and $s_m$ is the spread
of moment $m$ across resamples. The block length is one full session, longer
than the longest autocorrelation lag in the battery, so reshuffling blocks
preserves the dependence structure the ACF moments are meant to capture. The
$s_m$ are fixed for the entire calibration, which keeps the loss stationary and
$D(\theta)$ comparable across every evaluation. The interpretation is a
$z$-statistic: $|\widehat m_{\mathrm{sim}} - \widehat m_{\mathrm{emp}}|/s_m$ is
how many standard errors of pure historical sampling variation the simulated
moment misses by, and $D$ counts such standardised misses across the battery, so
a value of one is, loosely, a one-sigma miss per moment.

Equation~\eqref{eq:cal-loss} is the diagonal, absolute-deviation member of the
\citet{franke_westerhoff_2012} inverse-variability family, and the choice among
that family's members is deliberate. The full quadratic inverse-covariance form
$W=\widehat\Sigma^{-1}$ is ill-conditioned with this many moments and a single
historical sample. Inverse-\emph{variance} weighting, $1/s_m^2$, is
scale-pathological here: the minute sampling variance of the return standard
deviation would inflate its component by many orders of magnitude and collapse
the tail fit. Raw equal weights would let the large-scale moments bury the small
ones. Inverse-SD is the principled middle that makes the moments commensurate
while keeping $D$ stable and interpretable.

One property of \eqref{eq:cal-loss} is recorded for honest reporting. The KS
statistic is a two-sample distance whose sampling scale shrinks with the
\emph{simulated} sample size (it falls roughly as $1/\sqrt{n}$), so $s_{\mathrm{KS}}$
is computed at the simulation's own sample size rather than the full history's;
otherwise the KS component would be over-weighted and dominate the loss.
A consequence is that $D$ values are comparable only between runs with the same
number of simulated days and seeds, and all headline comparisons in the results
chapter therefore hold the seed count fixed.

\bigskip
\noindent
With the target and the loss fixed, the two methods that minimise $D(\theta)$ ---
an exhaustive grid on the true simulator and the surrogate-assisted search that
agrees with it --- and the validation protocol that guards the located optimum,
follow next.


% =====================================================================
% NOTES FOR THE AUTHOR
% ---------------------------------------------------------------------
% Citation keys used (all already resolve):
%   franke_westerhoff_2012, lamperti_2018, gao_2022_xgb, gao_2022_hfabm,
%   gao_2023, cont_2001, hill_1975, kunsch_1989, cont_stoikov_talreja_2008,
%   majewski_2018, simudyne_odd.
%
% OPTIONAL — to enrich the SMM lineage in §0 (only if you want the classical
% roots cited alongside Franke–Westerhoff). Not currently \cite'd above:
%
% @incollection{winker_gilli_maringer_2007,
%   author    = {Winker, Peter and Gilli, Manfred and Jeleskovic, Vahidin},
%   title     = {An objective function for simulation based inference on exchange rate data},
%   booktitle = {Journal of Economic Interaction and Coordination},
%   volume    = {2}, number = {2}, pages = {125--145}, year = {2007}}
%
% @article{grazzini_richiardi_2015,
%   author  = {Grazzini, Jakob and Richiardi, Matteo},
%   title   = {Estimation of ergodic agent-based models by simulated minimum distance},
%   journal = {Journal of Economic Dynamics and Control},
%   volume  = {51}, pages = {148--165}, year = {2015}}
% =====================================================================
